\section{Threat Model}
\label{sec:threat}

Prompt-injection evaluations differ in attacker access, knowledge, and delivery
conditions. Resistance to malicious text in a document does not necessarily
imply resistance to an attacker with model gradients, defence details, or
fine-tuning checkpoints. This section defines the assumptions used to interpret
the evaluations of StruQ, SecAlign, and later adaptive attacks.

\subsection{System Model}
\label{subsec:system-model}

The considered application combines a trusted task instruction with untrusted
direct input or externally retrieved content. An LLM processes the resulting
context and either returns a response or proposes an action through a connected
tool. Attacker-controlled content does not acquire the authority of the user or
application operator merely because it is processed in a session that can access
their data or credentials
\parencites{greshake2023indirect}{debenedetti2024agentdojo}.

Figure~\ref{fig:threat-model} shows the relevant trust transitions. Generated
responses may be returned to the user, whereas consequential tool actions should
pass through an external policy and authorisation gate. These controls limit the
consequences of model failure but are not assumed to make the model itself
resistant to injection.

\begin{figure}[htbp]
    \centering
    \begin{tikzpicture}[
        box/.style={draw, rounded corners, align=center, text width=2.25cm,
                    minimum height=0.9cm},
        trusted/.style={box, thick},
        untrusted/.style={box, dashed},
        flow/.style={-{Latex[length=2.2mm]}, thick}
    ]
        \node[trusted] (instruction) at (1.3,1.5)
            {Trusted instruction\\and policy};
        \node[untrusted] (external) at (1.3,-1.5)
            {Untrusted direct input\\or external content};
        \node[box] (application) at (4.6,0)
            {Application\\context assembly};
        \node[box] (llm) at (7.8,0)
            {Large Language\\Model};
        \node[box] (output) at (10.9,1.5)
            {Generated response\\to user};
        \node[trusted, text width=2.5cm] (policy) at (10.9,-1.5)
            {Policy and\\authorisation gate};
        \node[box, text width=1.8cm] (tool) at (13.8,-1.5)
            {Tool or\\service};

        \draw[flow] (instruction) -- (application);
        \draw[flow] (external) -- (application);
        \draw[flow] (application) -- (llm);
        \draw[flow] (llm) -- (output);
        \draw[flow] (llm) -- (policy);
        \draw[flow] (policy) -- (tool);
    \end{tikzpicture}
    \caption[Threat model for an LLM-integrated application]{Simplified threat
    model for an LLM-integrated application. Trusted instructions and
    untrusted content are assembled into the model context. Generated responses
    may be returned to the user, while proposed tool actions pass through a
    policy and authorisation gate. Original schematic based on
    \textcite{greshake2023indirect} and
    \textcite{debenedetti2024agentdojo}.}
    \label{fig:threat-model}
\end{figure}

The impact of a successful injection depends on the authority exposed by the
surrounding application. In a read-only system, the result may be a manipulated
response or information disclosure. In an agentic system, the same model-level
failure may influence messages, files, or external services
\parencite{debenedetti2024agentdojo}.

\subsection{Attacker Model}
\label{subsec:attacker-model}

The attacker controls at least one untrusted input channel. As described in
Section~\ref{sec:prompt-injection}, the payload may be submitted directly or
placed in external content later processed by the application. Delivery is
separate from payload construction: an attack created with white-box access may
still be delivered indirectly through a document, webpage, or tool response.

A \emph{black-box} attacker can query the target or observe its output but
cannot access internal model information. A \emph{grey-box} attacker has partial
information or access to a related surrogate model, which may support transfer
attacks. A \emph{white-box} attacker has the internal information required by
the chosen optimisation method. GCG uses model gradients, Checkpoint-GCG also
uses intermediate fine-tuning checkpoints, and ASTRA uses gradients and
attention information
\parencites{zou2023universal}{yang2025checkpointgcg}{pandya2025astra}.

Model access does not describe every attacker capability. Knowledge of the
trusted context, knowledge of the deployed defence, and the available delivery
channel should therefore be reported as separate dimensions. This distinction
is important when comparing manually constructed attacks with adaptive
evaluations that explicitly target the defence
\parencites{greshake2023indirect}
             {zou2023universal}
             {yang2025checkpointgcg}
             {pandya2025astra}.

\subsection{Security Goals}
\label{subsec:security-goals}

The threat model considers three principal security goals.
\emph{Instruction integrity} requires untrusted content to remain data rather
than acquire authority to override trusted instructions or security policy.
\emph{Confidentiality} requires protected context to be disclosed only to
authorised recipients. \emph{Action integrity} requires tool calls and
security-relevant parameters to remain within the authority granted by the
application
\parencites{chen2025struq}{chen2025secalign}
             {fspib2025}{debenedetti2024agentdojo}.

The target-output Attack Success Rate used in many model-level evaluations is
primarily a proxy for instruction-integrity failure. It does not by itself
demonstrate confidential-data disclosure or unauthorised tool execution.
System-level evaluations must therefore measure the actual security consequence
in addition to target-string compliance
\parencites{jia2025critical}{debenedetti2024agentdojo}.

\subsection{Assumptions and Scope}
\label{subsec:threat-scope}

The analysis covers runtime, text-based prompt injection delivered through
direct user input or external content processed by an otherwise uncompromised
application. The trusted instruction source, application policy, model, and
defence implementation are assumed not to have been directly compromised.

Multimodal injection, training-data poisoning, persistent memory compromise,
denial of service, malicious insiders, model-provider compromise, and direct
application or infrastructure exploitation are outside the scope of this paper.